% !PRINTED 505
\weberpage{505}
%<<CONT
\gesperrt{ist, in der die Coëfficienten $A_1, A_2, \ldots, A_m$ ganze rationale Functionale sind.}

Functionale, die nicht zu den ganzen gehören, werden wir gelegentlich auch der Kürze wegen als gebrochene Functionale bezeichnen.

Zur Rechtfertigung dieser Definition beweisen wir zunächst den Satz:

3. \gesperrt{Ist $\omega$ ein ganzes Functional nach der Definition 2. und zugleich rational, so ist es ein ganzes rationales Functional (nach der Definition 1.).}

Angenommen, es sei $\omega$ ein gebrochenes rationales Functional, und daher der absolute Werth von $\omega$ ein rationaler Bruch $p:q$, worin $p$ und $q$ ganze rationale Zahlen ohne gemeinsamen Theiler sind; dann ist $q\omega$ ein ganzes rationales Functional, dessen absoluter Werth $=p$, also relativ prim zu $q$ ist. Wenn aber andererseits $\omega$ zugleich ganz im Sinne der Definition 2. ist, so können wir
\[
\omega^m = -(A_1\omega^{m-1} + A_2\omega^{m-2} \ldots)
\]
setzen, woraus
\[
(q\omega)^m = -q[A_1(q\omega)^{m-1} + A_2\,q(q\omega)^{m-2} + \cdots) % [sic] opening square bracket is closed by a round parenthesis as printed
\]
folgt.

Hieraus aber ergiebt sich, dass der absolute Werth $p^m$ von $(q\omega)^m$ durch $q$ theilbar sein muss, was nur möglich ist, wenn $q=1$ ist. Die Definition 1. ist also in der Definition 2. als Specialfall enthalten.

Ist
\[
\Phi(t) = t^m + A_1t^{m-1} + A_2t^{m-2} + \cdots + A_m
\]
eine Function von $t$, in der die Coëfficienten $A_1, A_2, \ldots, A_m$ gebrochene rationale Functionale mit den Variablen $x, y, \ldots$, aber von der Variablen $t$ frei sind, so kann man eine Function $aE(x,y,\ldots)$ bestimmen, in der $a$ den Hauptnenner der absoluten Werthe von $A_1, A_2, \ldots$ und $E(x,y,\ldots)$ eine primitive ganze Function bedeutet, so dass
\weq{1}{aE(x\ y,\ldots)\,\Phi(t) = P(t,x,y,\ldots)} % [sic] printed without comma between x and y
selbst eine primitive ganze Function ist [§.\,136, 2.].

Zerfällt $\Phi(t)$ in zwei Factoren $\Phi_1(t)$, $\Phi_2(t)$ in $R$, ist also
\[
\Phi(t) = \Phi_1(t)\,\Phi_2(t),
\]
so bestimme man hiernach für die beiden Functionen $\Phi_1, \Phi_2$ die Factoren $a_1E_1, a_2E_2$, so dass
%CONT>>
